\documentclass[../main.tex]{subfiles}
\begin{document}
%==========================================TIÊU ĐỀ TẬP========================================
\begin{table}[h]
    \begin{adjustwidth}{-1cm}{}
        \begin{tabular}{>{\centering\arraybackslash}p{6cm}|>{\raggedright\arraybackslash}p{12.5cm}}
            \multirow{5}{6cm}
            % Cột bên trái: ÉPISODE và số tập
            {\rainbowcolor
                {\\[-40pt]
                    \flaregothic\fontsize{28pt}{28pt}\selectfont \centering
                    \color{eptype}LE PRÉLUDE\color{black}\\[12pt]
                    % \vnmdisp\fontsize{36pt}{36pt}\selectfont
                    % \color{epnum}NĂM
                    % \\[18pt]
                    % \flaregothic\fontsize{14pt}{14pt}\selectfont
                    % \color{epattr}BIENVENUE, \uline{\color{Banpire} Mel} !
                    % \flaregothic\fontsize{18pt}{18pt}\selectfont
                    % \color{epattr}ÉPISODE PERDU\\
                    \unrainbowcolor}
                \color{black}}
             &
            % Dòng thứ nhất: tên tập tiếng Nhật
            {\fotskip\fontsize{18pt}{18pt}\selectfont にじさんじの\ruby{世}{せ}\ruby{界}{かい}へようこそ！
            }      \\[6pt]
            % Dòng thứ hai: tên tập tiếng Anh
             & {\cafeteria\fontsize{18pt}{18pt}\selectfont Welcome to NIJISANJI World!}       \\[4pt]
            % Dòng thứ ba: tên tập tiếng Việt
             & {\vnmsans\fontsize{18pt}{18pt}\selectfont Thế giới của vùng đất Cầu Vồng}      \\[8pt]
            % Dòng thứ tư: Nhân vật xuất hiện
             & {\sen{\fontsize{14pt}{14pt}\selectfont Track used: Illusion (\ruby{幻}{まぼろし})}}
            \\[18pt]
            % % Dòng cuối cùng: Thời gian diễn ra của tập (thường sẽ là ngày phát hành)
            % & {\continuum\fontsize{16pt}{16pt}\selectfont Ngày phát hành: 12/09/2021}\\[18pt]
        \end{tabular}
    \end{adjustwidth}
\end{table}
%=========================================TRUYỆN CHÍNH========================================
\color{black}
\begin{center}
    \pov{Hikawa Kuon}
    \uline{Tiếp tục từ {\sen Partie 5}\\Vũ trụ 2434.}
\end{center}

\dots

\dots

``Ái cha cha\dots''

Tôi không biết mình đã ngất đi bao lâu.

Thứ đầu tiên tôi cảm nhận được là nền đất lạnh lẽo dưới lòng bàn tay. Một làn gió nhẹ lướt qua mặt, mang theo hương vị của khói xe, cát bụi\dots\ và thoảng trong đó là mùi gì đó quen lắm, như\dots\ bánh takoyaki mới nướng. Nhưng đầu óc tôi còn quá mụ mị để phân biệt được điều gì là thực, điều gì là ảo ảnh.

Một âm thanh mơ hồ vang lên từ xa: tiếng còi xe, tiếng bước chân, tiếng người gọi nhau trên phố\dots\ tất cả hòa trộn vào nhau như một bản hợp âm không đầu không cuối.

Tôi từ từ mở mắt ra.

Một bầu trời xanh thẳm hiện lên trên đầu, vắt ngang bởi những tòa nhà cao tầng phủ kính xanh. Mặt trời nhợt nhạt hắt xuống ánh sáng không gắt, đủ để khiến tôi nheo mắt, nhưng không quá chói để tôi phải quay đi.

Tôi ngồi dậy. Toàn thân ê ẩm như vừa trải qua một giấc mơ dài hàng thế kỷ. Không gian xung quanh tôi là một con phố rộng, lát nhựa sạch sẽ, nằm kẹp giữa hai dãy nhà cao tầng như một khe hẹp của đô thị hiện đại.

Nhưng có gì đó khiến tôi cảm thấy là lạ.

Tôi đã từng sống ở Tokyo (ảo và thực) phần lớn trong đời, trước khi quay về Kawachi khi tôi ngấp nghé bước sang tuổi mười bốn. Tôi ngồi dậy. Toàn thân ê ẩm như vừa trải qua một giấc mơ dài hàng thế kỷ. Không gian xung quanh tôi là một con phố rộng, lát nhựa sạch sẽ, nằm kẹp giữa hai dãy nhà cao tầng như một khe hẹp của đô thị hiện đại.

Nhưng\dots\ tôi lại cảm thấy có gì đó vừa lạ vừa quen ở nơi này.

Những bảng tên cửa hàng được viết bằng tiếng Nhật, đúng, nhưng phông chữ hơi lạ, logo hơi khác. Và kỳ lạ hơn, ở một tòa nhà lớn phía xa, tôi thấy một biểu tượng mà tôi chưa từng thấy: bốn ký tự N.J.S.J. được viết nghiêng, lồng ghép vào nhau như một thương hiệu.

Tôi không biết NJSJ là gì. Nhưng linh cảm bảo tôi rằng\dots\ nơi này không phải Nhật Bản mà tôi từng biết.

``Kaigou-chan? Suzuki?''

Tôi vội đảo mắt tìm quanh. Không thấy hai chị em đâu cả. Nỗi lo lắng dâng lên trong tôi như một làn khói đặc. Không lẽ nào họ lại lạc ở đâu đó rồi?

Nhưng ơn giời, chuyện đó không xảy ra. Phải mất gần hai mươi giây sau, tôi mới nghe thấy một tiếng rên khe khẽ ở bồn cây phía sau.

``\vie{Ugh\dots\ ai đẩy tôi xuống vậy trời\dots}'' giọng đó là của Kaigou-chan. Thật quen, đầy cáu kỉnh, y như mọi khi.

Tôi vội chạy lại. Cô nàng tóc xoăn cột đuôi ngựa đang lồm cồm ngồi dậy, ôm lấy vai.

``\vie{Sống rồi thì nói sớm, làm tôi cứ tưởng\dots} tôi chưa kịp nói hết câu thì lại nghe thấy tiếng hắt hơi nhỏ phía góc tường.

``\dots Onii-sama?''

Tôi quay đầu lại, thấy Suzuki đang dụi mắt, trông như một con mèo nhỏ vừa bị đánh thức khỏi giấc ngủ mùa đông. Cô bé vẫn mặc trang phục áo dài tay và váy dài như cũ, quần tất hơi rách ở đầu gối, bám đầy bụi đường.

``\vie{Em ổn chứ?}'' Tôi hỏi, rồi nhẹ tay phủi bụi trên vai em. Cô bé gật nhẹ, rồi dựa vào tay tôi, mắt vẫn lơ mơ.

``Khoan\dots\ còn một người nữa.'' Tôi lẩm bẩm. Chính xác hơn, một cá thể nữa đã đi theo chúng tôi trong cái lỗ xoáy kia.

``\eng{Tôi đây,}'' một giọng nói vang lên từ cổ tay tôi. Một dải sáng hình tam giác lập thể bật lên từ mặt đồng hồ. Giọng tiếng Anh lẫn Pháp khàn khàn ấy chẳng thể lẫn đi đâu được.

``Game?''

``\eng{Không nghĩ là tôi còn hoạt động được sau cú rơi đó đâu,}'' Game đáp. ``\eng{Nhưng rõ ràng, nơi này không phải nơi chúng ta đã ở.}''

``\eng{\dots Bằng cách nào ông chui tọt vào chiếc đồng hồ đeo tay của tôi vậy?}'' tôi nhướng mày, nhìn kỹ hơn vào chiếc đồng hồ.

``\eng{Tôi không biết. Lúc tôi tỉnh dậy thì đã thấy như thế này rồi,}'' ông thở dài. ``\eng{Cơ mà, dựa trên phân tích dựa trên tín hiệu GPS và kiến trúc khu vực. Mọi thứ đều khớp với Nhật Bản mà ta biết. Nhưng có lẽ đây không phải là nơi mà cậu từng biết.}''

Kaigou chống nạnh, cau mày nhìn quanh: ``\vie{Tôi thấy có gì đó sai sai. Kiểu, cứ như đây là nơi mà tôi từng biết, nhưng bị ai đó chỉnh sửa lại đôi chút.}''

Suzuki lẩm bẩm, nhìn vào một tấm áp phích của một tòa nhà chọc trời: ``\vie{Các bảng hiệu kia\dots\ hình như có tên của ai đó kìa.}''

Tôi liếc lên. Một biển hiệu quảng cáo cho món cơm cà ri in hình chibi một chàng trai đang nháy mắt: ``Collab đặc biệt với Mayuzumi-kun - chỉ tháng này!''

Tôi nhíu mày. ``Mayuzumi\dots\ Mayuzumi Kai\footnote{Sẽ xuất hiện trong {\flaregothic JOUR \hbrk 5}.}?''

Cô chị lớn liếc nhìn tôi: ``\vie{Người quen của cậu à?}''

``\vie{Không. Nhưng là một cái tên tôi từng nghe\dots\ từ Nijisanji.}''

Cô em nhỏ thì gật đầu liên tục: ``\vie{Đúng rồi ạ! Em nhớ chị em từng thấy người tên đó trong một clip collab. Còn kìa nữa\dots\ đó là cửa hàng ramen hợp tác với chị Hoshikawa Sara-san\footnote{Sẽ xuất hiện trong {\flaregothic JOUR \hbrk 7}.}\dots}''

Chúng tôi đứng im một lúc. Cả ba chúng tôi - cộng thêm Game - cùng nhìn nhau. Một cảm giác kỳ lạ len lỏi trong lòng tôi: như thể một thế giới khác đang mở ra ngay trước mắt. Thân thuộc, mà cũng hoàn toàn xa lạ.

``\vie{Tôi nghĩ\dots\ nơi này là thế giới của Nijisanji,}'' tôi nói, mắt vẫn còn trố ra vì ngạc nhiên.

``\vie{Có lẽ tôi đã nghe đến họ ở đâu đó rồi\dots} Kaigou đáp, xác nhận với bản thân.

``\vie{Em cũng có nghe đến, nhưng chưa từng nghĩ là\dots\ ta có thể rơi vào thế giới này}'' Suzuki thêm vào.

Game gật gù: ``\eng{Một chiều không gian song song. Có lẽ đây là một trong những lớp giao nhau giữa các dòng thực tại mà tôi từng giả định.}''

``\eng{Tôi không hiểu ông nói gì cả,}'' cô chih lớn thở dài.

Tôi cũng không hiểu nốt. Nhưng tôi biết một điều: chúng tôi không còn ở nơi cũ nữa, hay bất cứ nơi nào mà tôi hay hai chị em ở thế giới khác biết tới.

Chúng tôi bước đi, lạc lõng giữa thành phố rộng lớn. Người dân xung quanh không để ý tới ba kẻ ngoại lai đang bối rối trước mọi thứ. Tàu điện chạy êm ru trên cao, các quán ăn mở cửa, trẻ con nắm tay cha mẹ đi qua đường\dots\ Mọi thứ đều vận hành bình thường. Bình thường đến rợn người.

Cuối cùng, sau hơn nửa giờ lang thang vô định, chúng tôi dừng lại ở một quán ăn gia đình nhỏ góc phố. Ánh đèn vàng, cửa kính sạch bóng, bảng menu treo ngay lối vào: Hôm nay có combo đặc biệt kèm chữ ký từ cô nàng Gundoi-sensei\footnote{Nhân vật này không xuất hiện trong quyển sách. Tên đầy đủ: Gundou Mirei (\ruby{郡}{ぐん}\ruby{道}{どう}\ruby{美}{み}\ruby{玲}{れい}), và đã rời NIJISANJI vào 21/06/2023.}!

Tôi hít một hơi dài.

``\eng{Thôi thì\dots\ nếu chúng ta sẽ kẹt ở đây một thời gian, bắt đầu từ đâu đó ấm bụng trước đã.}''

Kaigou lườm tôi, nhưng cũng đẩy cửa bước vào. ``\eng{Đúng đó. Suốt mấy ngày nay lạc ở The Void, tôi đã có gì bảo vào bụng đâu.}''

Suzuki líu ríu theo sau hai chúng tôi, còn Game chỉ thở dài: ``\eng{Thế giới nào rồi cũng phải ăn, huh?}''

Cánh cửa đóng lại sau lưng chúng tôi.

Không gian bên trong quán ăn mang đậm nét của một quán ăn gia đình kiểu Nhật: ánh sáng vàng dịu hắt từ trần nhà, sàn gỗ sáng bóng, những dãy ghế nệm bọc da màu nâu nhạt chia thành từng khu riêng biệt. Cây xanh treo tường, tranh phong cảnh và sách bày gọn trên kệ. Một cảm giác ấm áp len vào lòng tôi, đối lập hẳn với sự lạnh lẽo vừa trải qua. Chúng tôi ngồi xuống một bàn gần cửa sổ, trong lúc tôi còn đang trấn tĩnh lại thì bắt đầu để ý tới xung quanh.

``Đây là\dots\ Nhật Bản à?'' Kaigou lẩm bẩm, chuyển về tiếng Nhật, đưa mắt đảo quanh quán. ``Giống mà không giống\dots''

``Cái mùi, ánh sáng, cả bảng thực đơn nữa. Đúng kiểu quán ăn gia đình Nhật,'' tôi nói, cũng không giấu được vẻ nghi hoặc. ``Nhưng vẫn có gì đó lạc lõng.''

Suzuki rụt vai, khẽ bám lấy tay áo tôi: ``Em thấy hơi run\dots\ Nhưng mà, nơi này\dots\ không có vẻ xấu nhỉ?''

``Ừ. Ít nhất thì\dots\ có người,'' Tôi nhìn quanh.

Khách trong quán khá đông. Và không bình thường chút nào.

Ở bàn ngay gần đó, một nhóm thanh niên đang nói chuyện rôm rả. Tôi liếc qua và nhận ra hai gương mặt quen thuộc từ ảnh quảng cáo lúc nãy. Nình như là\dots\ Hayato và Toya?

Ngồi ở gần đó một chàng trai có mái tóc trắng nhạt, làn da tái và cặp răng nanh rất rõ - chắc chắn là ma cà rồng - cùng một người khác trông như bước ra từ trò chơi nhập vai, mặc bộ áo giáp sáng loáng chẳng khác gì bản sao nam của Noel-chan.

Bên trong sâu hơn là một bàn bốn cô gái mặc đồng phục nữ sinh, đang vừa cười khúc khích vừa chia nhau dĩa bánh ngọt phủ kem. Một trong số đó khoác trên vai chiếc áo khoác gấu trúc to sụ, cái tai thú trên mũ cứ đung đưa theo từng cử động.

Kaigou nhìn thấy thì bật cười nhẹ: ``Cô bé kia mặc kiểu gì thế kia? Nhưng mà đáng yêu thật\dots''

``Em thấy như đang nhìn một nhóm nữ sinh trung học em thấy trên anime vậy,'' Suzuki nghiêng đầu.

Góc xa bên phải là một khung cảnh trầm hơn: ba người đang nhâm nhi trà nóng. Một cô gái mặc kimono chỉnh tề, mái tóc dài đen nhánh. Kế bên là một bà sơ đeo mạng che mặt, và cuối cùng là một cô gái với mái tóc đỏ rực và đôi mắt như lửa cháy, khí chất như một con rồng.

``\dots Có vẻ như mỗi bàn là một thế giới riêng vậy.'' Tôi khẽ nói.

Tôi quay đầu về phía bên trái, thấy một bàn nữa khá lạ: một người đàn ông trung niên mặc vest, trông phong trần; một nữ thần với mái tóc dài lấp lánh và trang phục đính đầy dải lụa uốn lượn; một cậu trai trẻ có vẻ ngang tuổi tôi với vẻ mặt trầm mặc, và một kẻ khác đội mũ ngược, áo khoác trễ, dáng điệu như một tên đầu gấu đang ngồi nghịch ly soda có mấy quả cherry nổi lềnh bềnh.

Và còn nữa. Rải rác khắp quán là những gương mặt không ai giống ai. Mỗi người đều như tỏa ra một bầu khí riêng biệt, có người thần thánh, có người yêu dị, có kẻ lại quá mức đời thường. Tôi đếm sơ sơ cũng phải hơn mười người tất cả.

``Nhìn kiểu gì cũng không phải khách bình thường,'' Kaigou nghiêng đầu. ``Chúng ta đang lạc vào cái gì thế này?''

``Đếm sơ sơ chắc cũng phải hơn chục người rồi,'' Tôi thở dài, tay chống cằm. ``Và nếu tất cả bọn họ đều như vậy thì\dots''

Suzuki nhìn quanh rồi lí nhí: ``Chuyến đi này có vẻ\dots\ không đơn giản rồi\dots''

``Đơn giản là cái gì?'' cô nàng cười mỉa. ``Thế giới của chị em mình chưa đủ phức tạp sao?''

``\textit{Chuyến đi này lành ít dữ nhiều đây\dots}'' Tôi lẩm bẩm, rồi thở dài khi nhìn đám đông trong quán, tay nắm chặt chiếc thìa cà-phê như để trấn an bản thân rằng mình vẫn còn thức.

Và thế là, mười hai ngày phiêu lưu của chúng tôi, ở thế giới của Nijisanji chính thức \emph{bắt đầu}.

\end{document}